\documentclass[11pt,a4paper]{article}

\usepackage[T1]{fontenc}
\usepackage{lmodern}
\usepackage[margin=20mm]{geometry}
\usepackage{amsmath,amssymb,amsthm}
\usepackage{microtype}

\theoremstyle{definition}
\newtheorem*{problem*}{Problem}

\setlength{\parindent}{0pt}
\setlength{\parskip}{5pt}

\begin{document}

\begin{center}
  {\Large\bfseries The sharp distance-to-circle decay rate for curve
diffusion - single number}
\end{center}

\section{Definitions}

Let $\Omega\subset\mathbb R^2$ be a smooth domain whose boundary
$\partial\Omega$ represents a fixed wall. For $t\geq0$, let
\begin{equation}
  \Gamma_i(t)=X_i([0,1],t),
  \qquad i=1,2,3,
\end{equation}
be three immersed curves. They meet at a common triple junction:
\begin{equation}
  X_1(0,t)=X_2(0,t)=X_3(0,t)=q(t),
\end{equation}
and their other endpoints lie on the wall:
\begin{equation}
  X_i(1,t)\in\partial\Omega,
  \qquad i=1,2,3.
\end{equation}

Orient each curve away from the triple junction. Let
\begin{equation}
  \tau_i,
  \qquad
  n_i,
  \qquad
  \kappa_i,
  \qquad
  V_i=\partial_tX_i\cdot n_i
\end{equation}
denote, respectively, the unit tangent, unit normal, scalar curvature, and
normal velocity of $\Gamma_i(t)$. Let $\sigma_i>0$ be the surface tensions
and $m_i>0$ the mobilities. Define the chemical potentials by
\begin{equation}
  \mu_i=\sigma_i\kappa_i.
\end{equation}
Here $\partial_{s_i}$ denotes differentiation with respect to arclength on
$\Gamma_i(t)$.

The first two curves evolve by surface diffusion, while the third curve
evolves by mean curvature flow:
\begin{equation}
  \begin{aligned}
    V_i
    &= -m_i\partial_{s_i}^2\mu_i,
    && i=1,2,
    && \text{on }\Gamma_i(t),\\
    V_3
    &= m_3\mu_3=m_3\sigma_3\kappa_3,
    &&
    && \text{on }\Gamma_3(t).
  \end{aligned}
  \label{eq:evolution}
\end{equation}

At the triple junction, impose the concurrency condition
\begin{equation}
  X_1(0,t)=X_2(0,t)=X_3(0,t)=q(t),
  \label{eq:concurrency}
\end{equation}
the kinematic compatibility conditions
\begin{equation}
  q_t(t)\cdot n_i(0,t)=V_i(0,t),
  \qquad i=1,2,3,
  \label{eq:kinematic}
\end{equation}
and Herring's law
\begin{equation}
  \sigma_1\tau_1(0,t)
  +\sigma_2\tau_2(0,t)
  +\sigma_3\tau_3(0,t)=0.
  \label{eq:herring}
\end{equation}
Assume in addition that
\begin{equation}
  \sigma_1=\sigma_2,
  \qquad
  2\sigma_1>\sigma_3.
\end{equation}

Continuity of chemical potential is imposed only between the two
surface-diffusion curves:
\begin{equation}
  \mu_1(0,t)=\mu_2(0,t).
  \label{eq:chemical-continuity}
\end{equation}
Mass-flux balance at the triple junction is likewise imposed only on these
two curves:
\begin{equation}
  m_1\partial_{s_1}\mu_1(0,t)
  +m_2\partial_{s_2}\mu_2(0,t)=0.
  \label{eq:flux-balance}
\end{equation}
Both arclength coordinates are oriented away from the triple junction. No
chemical-potential continuity or mass-flux condition is imposed on the
mean-curvature-flow branch $\Gamma_3(t)$.

Let $\nu_{\partial\Omega}$ denote the unit normal to the wall. All three
curves meet the wall at a right angle:
\begin{equation}
  n_i(1,t)\cdot
  \nu_{\partial\Omega}\bigl(X_i(1,t)\bigr)=0,
  \qquad i=1,2,3.
  \label{eq:contact-angle}
\end{equation}
The no-flux condition at the wall is imposed only on the two
surface-diffusion curves:
\begin{equation}
  \partial_{s_i}\mu_i(1,t)=0,
  \qquad i=1,2.
  \label{eq:wall-no-flux}
\end{equation}

The initial conditions are
\begin{equation}
  \Gamma_i(0)=\Gamma_i^0,
  \qquad i=1,2,3.
  \label{eq:initial-data}
\end{equation}
The initial network consists of immersed curves with a single common triple
junction. It satisfies Herring's law, the orthogonal wall-contact
conditions, and the chemical-potential and flux compatibility conditions
associated with the two surface-diffusion branches.

\section{Problem Formalization}

Initially parametrize each curve by its arclength coordinate $s_i$. A
counterexample purporting to disprove uniqueness solely through a choice of
parametrization is not admissible unless an explicit gauge map has first been
fixed. As part of the solution, develop a method for explicitly imposing such
a gauge map and use it in the analysis of the rest of the problem.

Prove the following statements.

\begin{enumerate}
  \item Establish short-time existence, uniqueness, continuous dependence on
  the initial data, and maximal $L^p$-regularity for the coupled system
  \eqref{eq:evolution}--\eqref{eq:initial-data}, for $p>5$ with the lower
  restriction on $p$ made as weak as possible, and under suitable initial
  compatibility conditions.

  More precisely, after expressing the evolving curves as normal graphs over
  a suitable reference network using the gauge developed above, prove that
  there exists $T>0$ such that
  \begin{equation}
    X_i
    \in
    W_p^1\bigl((0,T);L_p(0,1;\mathbb R^2)\bigr)
    \cap
    L_p\bigl((0,T);W_p^4(0,1;\mathbb R^2)\bigr),
    \qquad i=1,2,
  \end{equation}
  and
  \begin{equation}
    X_3
    \in
    W_p^1\bigl((0,T);L_p(0,1;\mathbb R^2)\bigr)
    \cap
    L_p\bigl((0,T);W_p^2(0,1;\mathbb R^2)\bigr).
  \end{equation}

  Show that the corresponding natural initial trace spaces are
  \begin{equation}
    X_i^0\in W_p^{4-4/p}(0,1;\mathbb R^2),
    \qquad i=1,2,
  \end{equation}
  and
  \begin{equation}
    X_3^0\in W_p^{2-2/p}(0,1;\mathbb R^2),
  \end{equation}
  subject to precisely those compatibility conditions whose traces are well
  defined in these spaces. In particular, verify the
  Lopatinskii--Shapiro complementing condition for the coupled junction and
  wall boundary operators.

  \item Find nontrivial assumptions on the initial network that guarantee
  global well-posedness.
\end{enumerate}


\end{document}